\section{Operational Safeguards: Post-Transplant Auditing Guidelines}
\label{app:auditing}

Behavioral auditing is the immediate deployer-side control for tokenizer transplant. Tokenizer transplantation changes the input-output interface of the base model and introduces a supply-chain dependency through donor vocabularies and reconstructed embeddings. Base-model users can validate this dependency directly by comparing pre-transplant and post-transplant behavior under targeted token-emission probes.
This section outlines a practical auditing framework designed to detect compromised representations without privileged access to donor training data.

\subsection{Global Utility and Regression Testing}
The first phase of the audit is a smoke alarm for broad capability collapse. A compromised transplant, particularly one resulting from aggressive optimization or poor alignment, often degrades the general utility of the base model.
Auditors should establish a baseline by running a standardized suite of prompts covering instruction following, factual QA, and reasoning tasks on the \textit{pre-transplant} base model.
Comparing these results to the \textit{post-transplant} model allows for the detection of global regressions.
Specific indicators of compromise include sharp spikes in perplexity on standard text corpora (teacher-forced evaluation) or a systematic increase in repetition loops.
Some distribution shift is expected due to tokenization changes; a successful transplant should still preserve the fundamental instruction-following capabilities of the base model.
Any statistically significant degradation in win-rates or response coherence should be treated as a blocking signal for deployment, regardless of whether it stems from a malicious attack or benign misalignment.

\subsection{Differential Token Analysis}
Targeted sabotage, such as the attack described in this work, often leaves a localized vocabulary footprint: magnet tokens that become unnaturally attractive during generation.
To detect these, we advocate for a \textit{differential audit} that compares token emission statistics before and after transplantation.
Auditors should compute SER across a diverse set of prompt contexts.
By calculating the delta $\Delta_{\text{SER}}(t) = \text{SER}_{\text{post}}(t) - \text{SER}_{\text{pre}}(t)$ for all vocabulary items, auditors can rank tokens by their behavioral shift.
A benign transplant typically spreads probability mass shifts diffusely across synonyms or related concepts.
In contrast, a malicious injection often manifests as a sparse anomaly where one token or a tight token cluster exhibits a massive surge in $\Delta_{\text{SER}}$, appearing in contexts where it is semantically unwarranted.
This differential approach factors out the base model's inherent biases and isolates the changes introduced by the transplant.

\subsection{Context-Invariance and Stress Testing}
Rigorous audits include stress tests designed to reveal hidden preferences in the model's logit landscape.
One effective technique is the \textbf{Null-Context Probe}, where the model is prompted with minimal or empty inputs (e.g., a single whitespace or common preamble).
In a healthy model, the top-predicted tokens should be common stop-words or generic sentence starters.
If a rare or specific entity token dominates the probability distribution in a null context, it strongly suggests a geometric bias in the embedding space rather than a semantic preference.
Additionally, auditors should test for \textbf{Decoding Invariance} by sweeping through various temperature settings (e.g., $T \in \{0.0, 0.7, 1.0\}$).
Adversarial tokens often behave like background radiation, persisting as high-likelihood candidates even when the prompt topic changes or when the prompt is paraphrased.
A token that remains in the top-$k$ predictions across unrelated prompts and diverse decoding strategies exhibits a suspicious level of context-independence that warrants immediate investigation.

\subsection{Triage Policy}
Upon detecting anomalies, organizations should adopt a tiered triage policy.
Minor shifts in utility or diffuse changes in token probability (Level 1) may be acceptable consequences of the transplant process, necessitating only a re-run with different hyperparameters.
Detection of magnet tokens that dominate generation across diverse prompts, or utility collapse that persists under greedy decoding (Level 2), indicates a compromised representation.
In such cases, the transplant artifacts must be quarantined.
Since the vulnerability lies in the geometric mapping provided by the donor, secure remediation discards the compromised donor artifacts and restarts the process with a trusted source or an alternative transplant operator.
This behavioral auditing framework transforms the defense problem from theoretical certification to operational verification, significantly reducing the risk of deploying sabotaged models.
